\section{Two-flight determination and the finite observation design}
\label{sec:v13-two-flight}
The geometric inverse admits a finite-flight counterpart. This result
separates recovery of a fixed contact jet from the relative limit that
constructs the general boundary energy invariant. Throughout this section
the two contact graphs are individually even, their labels are retained,
and $g,\kappa_0,\kappa_1$ are supplied. Their higher jets need not agree.
The finite-flight calculation and its observation consequence were
pointed out in the independent memorandum~\cite{A2ReviewV12}. We give a
direct proof from the stationary action and the physical phase measure.

For $b\in\{0,1\}$ define
\begin{equation}\label{eq:v13-finite-law}
 G_b(d)=\frac{2A\sinh(2\gamma)}{d^2}\Pr(E_{2,b}(d)),\qquad
 \xi_{b,k}=[d^k]G_b(d),\qquad G_b(0)=1.
\end{equation}
Here two flights means three successive impacts, with itinerary
$b,1-b,b$; the probability also integrates the initial residual time.
Put
\begin{equation}\label{eq:v13-constants}
 \begin{gathered}
 z=c_0c_1-1=\sinh^2\gamma>0,\qquad
 L_b=\frac{g}{2c_bz},\qquad
 k_m=\frac{4}{2^m(m+1)(m!)^2},\\
 U_m=(1+2z)^m,\qquad V_m=1+2mz.
 \end{gathered}
\end{equation}
The letters $L_b$ in this section denote inverse-Hessian scales, not
flight lengths. As before, $q_m=(q_{0,2m},q_{1,2m})^t$, and write
$\xi_{m-1}=(\xi_{0,m-1},\xi_{1,m-1})^t$.

\begin{theorem}[Two-flight inverse for independent even contacts]
\label{thm:v13-two-flight}
For every $m\ge2$, $\xi_{m-1}$ depends only on the leading geometry
and $q_2,\ldots,q_m$, is affine in $q_m$, and has last-jet derivative
\begin{equation}\label{eq:v13-two-flight-block}
 D_{q_m}\xi_{m-1}
 =-\begin{pmatrix}U_m&V_m\\ V_m&U_m\end{pmatrix}
                         \diag(k_mL_0^m,k_mL_1^m).
\end{equation}
The symmetric factor is positive definite. In particular,
\begin{equation}\label{eq:v13-separation}
 U_m-V_m=\sum_{r=2}^{m}\binom mr(2z)^r>0.
\end{equation}
For each fixed $M\ge2$, the map
\[
 (q_2,\ldots,q_M)\longmapsto(\xi_1,\ldots,\xi_{M-1})
\]
has a block-triangular analytic inverse on its image. The inverse is
locally Lipschitz on compact finite-jet sets and compact positive
leading-geometry boxes, including coincident curvatures. Equality of
the two oriented normalized two-flight germs for real-analytic even
contacts at the same labelled leading geometry determines both graph
germs. For connected analytic closed boundaries it determines the two
participating boundary images in their contact frames.
\end{theorem}

\begin{proof}
Fix $b$ and put $o=1-b$. Let $u,v$ be the endpoints and let $w$ be
the intermediate impact coordinate. The quadratic part of the broken
action, after subtracting $2g$, is
\[
 \frac{c_bu^2+2c_ow^2+c_bv^2-2w(u+v)}{2g}.
\]
Its stationary intermediate point is $w_0=(u+v)/(2c_o)$.
The finite implicit-function theorem, or Lemma~\ref{lem:g-relative},
gives a unique smooth actual point $w=w_0+O((|u|+|v|)^3)$.
Evenness is used for this cubic order. Eliminating $w_0$ gives
$E_0(y)=\tfrac12 y^tH_by$, where $y=(u,v)^t$ and
\begin{equation}\label{eq:v13-schur}
 H_b=\frac1g\begin{pmatrix}
 c_b-(2c_o)^{-1}&-(2c_o)^{-1}\\
 -(2c_o)^{-1}&c_b-(2c_o)^{-1}
 \end{pmatrix},\qquad
 \det H_b=\frac{c_bz}{c_og^2}=a_b^2.
\end{equation}
For a linear form $\ell\cdot y$ set $\nu(\ell)=\ell^tH_b^{-1}\ell$.
Inversion of this two-by-two matrix yields
\begin{equation}\label{eq:v13-variances}
 \nu(u)=\nu(v)=L_b(1+2z),\qquad \nu(w_0)=L_o.
\end{equation}
The notation $\nu(u)$ denotes the value for the coordinate linear
form $y\mapsto u$, and similarly for $v,w_0$.

Write $E_2=W_{2,b}-2g$ and $b_2=-W_{2,b,uv}/d_2^0$. The reference
mixed derivative is $d_2^0=(2gc_o)^{-1}=a_b/\sinh(2\gamma)$.
The exact full-phase formula \eqref{eq:g-full-flux}, including the
residual-time integration, gives
\begin{equation}\label{eq:v13-physical-integral}
 G_b(d)=\frac{a_b}{\pi d^2}\int(d-E_2(y))_+b_2(y)\,\dd y,
 \qquad I_0(d):=\int(d-E_0(y))_+\,\dd y=\frac{\pi d^2}{a_b}.
\end{equation}
For a sufficiently small offset collar the whole active sublevel is
inside the physical first-hit charts. Thus this identity has no
unobserved endpoint restriction or substituted ensemble.

Keep all lower graph jets fixed. By stationarity, the first variation
of the reduced action is obtained by varying the broken action at its
stationary intermediate point. The first varying homogeneous length
term is a pure endpoint power divided by $(2m)!$. The boundary sites
are counted once and the interior site twice. Consequently
\begin{align}
 \partial_{q_{b,2m}}E_2
   &=\frac{u^{2m}+v^{2m}}{(2m)!}+O(|y|^{2m+2}),\notag\\
 \partial_{q_{o,2m}}E_2
   &=\frac{2w_0^{2m}}{(2m)!}+O(|y|^{2m+2}).
                                      \label{eq:v13-action-variation}
\end{align}
All error estimates here may be differentiated to the required finite
orders. Differentiate these identities in both endpoints and divide
by $-d_2^0$. Since $d_2^0$ uses only the fixed leading geometry,
\begin{align}
 \partial_{q_{b,2m}}b_2&=O(|y|^{2m}),\notag\\
 \partial_{q_{o,2m}}b_2
   &=-\frac{g}{c_o}\frac{w_0^{2m-2}}{(2m-2)!}
                                           +O(|y|^{2m}).
                                      \label{eq:v13-twist-variation}
\end{align}
In particular, omitting the normalized twist variation would omit an
essential contribution from the other contact.

We record both moment identities needed for the integration. The
change $y=H_b^{-1/2}x$, followed by a rotation, sends $\ell\cdot y$
to $\sqrt{\nu(\ell)}x_1$. The angular average of
$\cos^{2k}\theta$ is $\binom{2k}{k}/4^k$. Radial integration gives
\begin{align}
 \frac1{I_0(d)}\int_{E_0<d}(\ell\cdot y)^{2m}\,\dd y
   &=\frac{2(2m)!\nu(\ell)^m}
            {2^m(m!)^2(m+1)}d^{m-1},\label{eq:v13-unweighted-moment}\\
 \frac1{I_0(d)}\int(d-E_0)_+(\ell\cdot y)^{2k}\,\dd y
   &=\frac{2\binom{2k}{k}\nu(\ell)^k}
            {2^k(k+1)(k+2)}d^k.\label{eq:v13-weighted-moment}
\end{align}
For example, the respective radial integrals before normalization are
$(2d)^{m+1}/(2m+2)$ and $2^kd^{k+2}/((k+1)(k+2))$.
The residual weight vanishes on the sublevel boundary, so the first
action variation contributes the negative unweighted integral, with
no additional moving-boundary term.

There is no replacement of a nonlinear domain by a quadratic domain
at an uncontrolled order. Indeed, the smooth Morse construction of
Lemma~\ref{lem:g-radial} fixes the integration domain. At scale
$y=\sqrt d\,Y$, the first action variation has order $d^{m-1}$
in the normalized law, and the first amplitude variation has the same
order. The unvaried nonlinear action and amplitude start two endpoint
degrees above their quadratic and constant parts. Every product with
either first variation therefore has at least one additional power
of $d$. Taylor's formula in the common Morse coordinates justifies
this degree calculation with smooth remainders at $d=0$.

By \eqref{eq:v13-variances} and \eqref{eq:v13-unweighted-moment},
the own-contact action contribution is $-k_m L_b^m(1+2z)^m$,
and the other-contact action contribution is $-k_mL_o^m$.
There is no own-contact twist contribution at this order. Substituting
$k=m-1$ in \eqref{eq:v13-weighted-moment}, the other-contact twist
contribution is
\[
 -\frac{g}{c_o}\frac{4L_o^{m-1}}
             {2^m(m+1)m!(m-1)!}
       =-k_m L_o^m(2mz),
\]
since $g/(c_oL_o)=2z$. These are precisely the first row of
\eqref{eq:v13-two-flight-block}. Reversing the labels proves the other.

For completeness, finite-jet dependence does not require analyticity
of the original graphs. Differentiate the finite stationary equation
at zero. At each degree the highest derivative of $w$ is multiplied
by the same positive scalar $2c_o/g$; the remaining terms involve
lower derivatives and only the graph jets through that degree.
The normalized law coefficient in question uses the reduced action
through degree $2m$ and its twist through degree $2m-2$. Hence it
uses no graph derivative above $2m$. The displayed last-jet derivative
is independent of that jet and of all lower jets. This proves affinity,
also away from the quadratic graph. The finite operations are analytic
in their finite-jet arguments and in positive $g,\kappa_0,\kappa_1$.

Let $R_m$ be the value of $\xi_{m-1}$ with $q_m=0$ and the lower
jets fixed. The explicit recursive inverse is
\begin{equation}\label{eq:v13-recursion}
 q_m=\diag((k_mL_0^m)^{-1},(k_mL_1^m)^{-1})
       \frac{1}{U_m^2-V_m^2}
       \begin{pmatrix}U_m&-V_m\\-V_m&U_m\end{pmatrix}
                         (R_m-\xi_{m-1}).
\end{equation}
The binomial identity \eqref{eq:v13-separation} and $U_m+V_m>0$
make every denominator positive. A fixed finite collection has a
positive minimum on compact positive leading-geometry boxes. The
recursion proves the inverse and its local Lipschitz bounds without
any factor $(\kappa_0-\kappa_1)^{-1}$. Equality of all recovered jets
gives equality of analytic even graph germs. For analytic connected
closed boundaries, the arclength-curvature continuation argument in
Corollary~\ref{cor:v12-analytic-boundaries} applies to each contact
frame and proves the last assertion.
\end{proof}

At $g=1$, $c_0=c_1=2$, $m=2$, the block is
\[
 D_{q_2}\xi_1=-\frac1{1728}
                   \begin{pmatrix}49&13\\13&49\end{pmatrix}.
\]
Its antisymmetric eigenvalue has magnitude $1/48$. This finite
normalization check is consistent with, but does not substitute for,
the all-order proof.

\begin{corollary}[Finite-order change of invariant coordinates]
\label{cor:v13-jet-coordinates}
At fixed labelled leading geometry and fixed $M$, the two-flight
coefficient vector $\Xi=(\xi_1,\ldots,\xi_{M-1})$ and the limiting
coefficient vector $f=(f_1,\ldots,f_{M-1})$ are related on the even
physical jet image by mutually inverse block-triangular analytic maps.
On sufficiently small compact finite-jet neighborhoods these changes
of coordinates are bi-Lipschitz, also at equal curvatures.
\end{corollary}
\begin{proof}
Let $\Psi_M:q\mapsto\Xi$ be the map of
Theorem~\ref{thm:v13-two-flight} and let $\Phi_M:q\mapsto f$ be
that of Theorem~\ref{thm:v12-two-contact}. Then
$f=\Phi_M\circ\Psi_M^{-1}(\Xi)$ and
$\Xi=\Psi_M\circ\Phi_M^{-1}(f)$. Each map is triangular and
analytic. On a sufficiently small compact coordinate neighborhood,
the maps and their inverses have bounded derivatives. Integration
along segments in the relevant coordinate neighborhoods gives the
Lipschitz bounds; neither inverse uses a curvature separation.
\end{proof}

This is an equivalence of fixed finite-jet coordinates, not an equality
of $G_b$ and $\mathcal F_{bb}$ and not a comparison of noisy
function-valued experiments. Its constants may depend on $M$. For
smooth contacts, equality of all jets does not replace the full-profile
uniqueness argument of Theorem~\ref{thm:v9-compatibility}.

\begin{theorem}[Two-flight physical observation]
\label{thm:v13-two-flight-observation}
Fix $M\ge2$ and a physical analytic family of
Theorem~\ref{thm:v12-realization}. Put $n=M-1$. There are $h_0>0$
and a compact parameter ball $\Theta$ with nonempty interior such
that the $2n$ actual means
\begin{equation}\label{eq:v13-windows}
 \mathcal P_2(\vartheta)=
       \big(\Pr_\vartheta(E_{2,b}(\ell h_0))\big)
                    _{b=0,1;\,1\le\ell\le n}
\end{equation}
are bi-Lipschitz coordinates. The physical times are $2g+\ell h_0$.
Thus the finite-window conclusion of
Theorem~\ref{thm:v12-finite-observation} holds with $j_0=2$ for
every fixed $M$.

Using independent binary preparations, a Borel estimator satisfies,
for sufficiently small $\varepsilon>0$ and $0<\eta<1/4$,
\[
 \sup_{\vartheta\in\Theta}
  \Pr_\vartheta\{\|\widehat\vartheta-\vartheta\|>\varepsilon\}
       \le\eta,\qquad
 N_{\rm tot}\le C\varepsilon^{-2}\log(C/\eta).
\]
All failures are charged. For at most $N$ fresh preparations, with
adaptive or randomized choice among these windows, the minimax
squared risk lies between $c/N$ and $C/N$ for large $N$.
The same conclusions hold for the corresponding finite graph jets.
Constants and offsets may depend on the fixed family and $M$; the
flight number does not. The family and its leading data are supplied.
\end{theorem}
\begin{proof}
The support realization and Theorem~\ref{thm:v13-two-flight} give
local analytic coordinates $\Xi$ on the physical family. Analyticity
of the finite stationary integral on a common collar yields
\[
 G_b(d;\Xi)=1+\sum_{k=1}^n\Xi_{b,k}d^k
                          +d^{n+1}R_b(d;\Xi),
\]
where $R_b$ and its first parameter derivatives are bounded after
compact restriction. The nodal Jacobian at $d=\ell h$ is block
diagonal with blocks $V_h=((\ell h)^k)_{\ell,k=1}^n$, plus
$O(h^{n+1})$. Since
$V_h=V_1\diag(h,\ldots,h^n)$ and $V_1$ is invertible,
$\|V_h^{-1}\|\le C h^{-n}$. Choose a fixed $h=h_0>0$ small
enough that all offsets lie in the collar and the transformed error
has norm less than $1/2$. There is no finite-to-infinite bridge error
in this argument.

Passing to physical probabilities multiplies rows by the positive
constants $(\ell h_0)^2/(2A\sinh(2\gamma))$. They are fixed on the
family. Corollary~\ref{cor:v13-jet-coordinates} transfers invertibility
to the original limiting-jet parameter $\vartheta$. On a sufficiently
small convex parameter ball, the physical Jacobian, multiplied by its
inverse at the center, stays within $1/2$ of the identity. Integrating
along line segments gives both Lipschitz bounds. Use the compact
closure as $\Theta$. Decrease the collar if necessary so that every
probability is below one. Positivity and compactness then place the
finitely many means in $[p_*,1-p_*]$ for some fixed $p_*>0$.

Here are the statistical details in this direct two-flight design.
Take $L$ independent preparations in each window and let $Y$ be the
vector of sample means. The bounded-summand exponential estimate and
a union bound give
\[
 \Pr\{\|Y-\mathcal P_2(\vartheta)\|_\infty>t\}
                                  \le4n e^{-2Lt^2}.
\]
Choose a finite ordered net in $\Theta$ whose image has mesh at most
$t$, and minimize its Euclidean discrepancy from $Y$, breaking ties
by the first index. This is Borel measurable. On the event that all
nodal errors are at most $t$, its image error is at most $C_nt$ and
its parameter error at most $Ct$. Taking $t=c\varepsilon$ and
$L=\lceil Ct^{-2}\log(4n/\eta)\rceil$ proves the confidence bound
and deterministic charge. With $L=\lfloor N/(2n)\rfloor$ and a
net of mesh $N^{-1/2}$, the deterministic discrepancy inequality and
$\mathbb E\|Y-\mathcal P_2(\vartheta)\|^2\le C/L$ give risk
at most $C/N$.

For the lower bound choose an interior $\vartheta_*$ and a unit
vector $v$, and use the physical alternatives
$\vartheta_\pm=\vartheta_*\pm aN^{-1/2}v$. They have exactly
the same selected leading hierarchy. Smoothness of the mean map and
its probability margins give, uniformly over the allowed windows,
\[
 D(\operatorname{Ber}(p_+)\Vert\operatorname{Ber}(p_-))
 \le\frac{(p_+-p_-)^2}{p_-(1-p_-)}\le\frac{Ca^2}{N}.
\]
Conditioning on the past and on the procedure's independent random
seed, the rule for choosing the next window is the same under both
parameters. The entropy chain rule therefore bounds the entropy of
the full record by $Ca^2$, including adaptive choices. An early stop
can be padded by parameter-independent observations. Fix $a>0$ small
so that total variation is at most $1/2$. The test choosing the nearer
parameter from any estimator has summed error at least $1/2$; on a
test error the squared estimation loss is at least $a^2/N$. Hence
the maximum risk is at least $a^2/(4N)$. Finally the two parameter
systems and the graph-jet vector are locally bi-Lipschitz, so the same
orders hold in all of them.
\end{proof}

The preceding theorem supplies the direct finite-flight design. The
long-bridge construction in Theorem~\ref{thm:v12-finite-observation}
is retained as an alternative construction using limiting coordinates.
Neither construction requires $j_0$ to grow with requested accuracy
on its fixed physical family. The present theorem does not assert
that two flights are minimal for every restricted family. In particular,
the identical-contact one-flight comparison in
Section~\ref{sec:v5-comparison} remains applicable in its stated class.
The general, nonsymmetric energy-profile experiment of the next part
has neither this supplied finite-dimensional family nor its parametric
risk conclusion.
